% 

\documentclass[twoside]{article}
\setlength{\oddsidemargin}{0.25 in}
\setlength{\evensidemargin}{-0.25 in}
\setlength{\topmargin}{-0.6 in}
\setlength{\textwidth}{6.5 in}
\setlength{\textheight}{8.5 in}
\setlength{\headsep}{0.75 in}
\setlength{\parindent}{0 in}
\setlength{\parskip}{0.1 in}

%
% ADD PACKAGES here:
%

\usepackage{amsmath,amsfonts,amssymb,array,graphicx,arydshln}


\newcounter{lecnum}
\renewcommand{\thepage}{\thelecnum-\arabic{page}}
\renewcommand{\thesection}{\thelecnum.\arabic{section}}
\renewcommand{\theequation}{\thelecnum.\arabic{equation}}
\renewcommand{\thefigure}{\thelecnum.\arabic{figure}}
\renewcommand{\thetable}{\thelecnum.\arabic{table}}

%
% The following macro is used to generate the header.
%
\newcommand{\lecture}[4]{
   \pagestyle{myheadings}
   \thispagestyle{plain}
   \newpage
   \setcounter{lecnum}{#1}
   \setcounter{page}{1}
   \noindent
   \begin{center}
   \framebox{
      \vbox{\vspace{2mm}
    \hbox to 6.28in { {\bf EE502 - Linear Systems Theory II
	\hfill Spring 2023} }
       \vspace{4mm}
       \hbox to 6.28in { {\Large \hfill Lecture #1 \hfill} }
       \vspace{2mm}
       \hbox to 6.28in { {\it Lecturer: #2 \hfill } }
      \vspace{2mm}}
   }
   \end{center}
   \markboth{Lecture #1}{Lecture #1}

   \vspace*{4mm}
}

\renewcommand{\cite}[1]{[#1]}
\def\beginrefs{\begin{list}%
        {[\arabic{equation}]}{\usecounter{equation}
         \setlength{\leftmargin}{2.0truecm}\setlength{\labelsep}{0.4truecm}%
         \setlength{\labelwidth}{1.6truecm}}}
\def\endrefs{\end{list}}
\def\bibentry#1{\item[\hbox{[#1]}]}


\newcommand{\fig}[3]{
			\vspace{#2}
			\begin{center}
			Figure \thelecnum.#1:~#3
			\end{center}
	}

% Use these for theorems, lemmas, proofs, etc.
\newtheorem{theorem}{Theorem}[lecnum]
\newtheorem{lemma}[theorem]{Lemma}
\newtheorem{proposition}[theorem]{Proposition}
\newtheorem{claim}[theorem]{Claim}
\newtheorem{corollary}[theorem]{Corollary}
\newtheorem{definition}[theorem]{Definition}
\newenvironment{proof}{{\bf Proof:}}{\hfill\rule{2mm}{2mm}}
\newtheorem{exmp}[theorem]{Ex}

% **** IF YOU WANT TO DEFINE ADDITIONAL MACROS FOR YOURSELF, PUT THEM HERE:

\begin{document}

% Lecture Details
\lecture{11}{Asst. Prof. M. Mert Ankarali}


%%%%%%%%%%%%%%%%%%%%%%%%%%

\section{Observability in CT-LTI Systems}

In the context of observability of dynamical systems, it 
turns out that it is more convenient to think in terms of
``unobservable states'' and then connect it to the 
concept of observability and fully observable systems,
as reflected in the following definitions.

\textbf{Definition:} For an LTI continuous-time state-space representation
%
\begin{align*}
  \dot{x} &= A x + B u
\\
  y &= C x + D u
\end{align*}
%
A state $x_u$ is said to be unobservable over $t \in [0 , T )$,
if with $x(0) = x_u$ and $\forall \, u(t)$ we get the same $y(t)$ as we would with $x(0) = 0$. 

The set, $\bar{\mathcal{O}}_T$, of all unobservable states over $t \in [0 , T )$ forms a vector space, $\bar{\mathcal{O}}_T \subset 
\mathbb{R}^n$, and the system is called fully observable over
$t \in [0 , T )$, if $\mathrm{dim}[\bar{\mathcal{O}}_T] = 0$.

Note that for linear dynamical systems, the observability of a state and
system is independent of $u(t)$; in that respect, we will 
only analyze the zero-input response of the system in our derivations.

\textbf{Theorem:} The following statements regarding an initial state $x_u$ are equivalent:
\begin{enumerate}
    \item $x_u \in \bar{\mathcal{O}}_T$ (Unobservable over $t \in [0, T)$)
    \item $x_u \in \bar{\mathcal{O}}_{\tau}$ for all $\tau > 0$
    \item $x_u \in \mathcal{N}[ \mathbf{O} ]$, where the observability matrix is
    \begin{align*}
      \mathbf{O} = \begin{bmatrix}
        C \\ C A \\ C A^2 \\ \vdots \\ C A^{n-1}
      \end{bmatrix}
    \end{align*}
\end{enumerate}

\textbf{Proof:} \\
\textit{Step 1: (3) $\implies$ (2).}\\ 
Let $x_u \in \mathcal{N}[ \mathbf{O} ]$, then  
%
\begin{align*}
  \mathbf{O} x_u = 0 \implies C x_u = 0, \ C A x_u = 0, \ \dots \ , C A^{n-1} x_u = 0
\end{align*}
%
By the Cayley-Hamilton theorem, we can conclude that $C A^{l} x_u = 0$ for all integers $l \ge 0$. 
Evaluating the zero-input output response of the system with $x(0) = x_u$:
%
\begin{align*}
  y_{zi}(\tau) = C e^{A \tau} x_u = \sum_{l=0}^{\infty} \frac{\tau^l}{l!} C A^l x_u = 0
\end{align*}
%
This holds for all $\tau > 0$, implying $x_u \in \bar{\mathcal{O}}_{\tau}$.

\textit{Step 2: (2) $\implies$ (1).}\\ 
This is trivially true. If the state is unobservable for all $\tau > 0$, it is certainly unobservable for any specific $T > 0$.

\textit{Step 3: (1) $\implies$ (3).}\\ 
If $x_u \in \bar{\mathcal{O}}_T$, then the zero-input output response is identically zero over the interval:
%
\begin{align*}
  y_{zi}(t) = C e^{A t} x_u = 0 , \quad \forall t \in [0,T]
\end{align*}
Taking successive derivatives with respect to $t$ and evaluating at $t=0$ yields:
\begin{align*} 
  y_{zi}(0) &= 0 \implies C x_u = 0 \\
  \left[\frac{d}{dt}y_{zi}(t) \right]_{t=0} &= 0 \implies C A x_u = 0 \\
  &\vdots \\
  \left[\frac{d^{n-1}}{dt^{n-1}}y_{zi}(t) \right]_{t=0} &= 0 \implies C A^{n-1} x_u = 0
\end{align*}
This means $\mathbf{O} x_u = 0$, or equivalently, $x_u \in \mathcal{N}[ \mathbf{O} ]$.

Similar to reachability, we show that for CT-LTI systems,
observability and the unobservable (and observable) subspace
are independent of time. 

\begin{exmp}
\textbf{Problem:} Consider the following sparse $3 \times 3$ system:
\begin{align*}
    A = \begin{bmatrix} 0 & 1 & 0 \\ 0 & 0 & 1 \\ 0 & 2 & -1 \end{bmatrix} \quad , \quad C = \begin{bmatrix} 0 & 1 & 0 \end{bmatrix}
\end{align*}
Determine if the system is fully observable. If not, find the unobservable subspace.

\textbf{Solution:}
Let's construct the observability matrix $\mathbf{O}$ to check the observability rank condition:
\begin{align*}
    \mathbf{O} = \begin{bmatrix} C \\ C A \\ C A^2 \end{bmatrix} = \begin{bmatrix} 0 & 1 & 0 \\ 0 & 0 & 1 \\ 0 & 2 & -1 \end{bmatrix}
\end{align*}
We can observe that the first column is entirely zero, while the second and third columns are linearly independent. Thus, $\mathrm{rank}(\mathbf{O}) = 2 < 3$, meaning the system is not fully observable. The unobservable subspace is simply the null space of $\mathbf{O}$:
\begin{align*}
    \bar{\mathcal{O}} = \mathcal{N}[\mathbf{O}] = \mathrm{span} \left\{ \begin{bmatrix} 1 \\ 0 \\ 0 \end{bmatrix} \right\}
\end{align*}
This tells us that the initial condition of the first state variable, $x_1(0)$, is completely hidden and has absolutely no effect on the output $y(t)$ for any $t \ge 0$.
\end{exmp}

\subsection{CT Observability Grammian}

For a CT-LTI system, the observability Grammian over the interval $[0, t]$ is defined as
%
\begin{align*}
    \mathbf{Q}(t) = \int\limits_{0}^{t} e^{A^T \tau} C^T C e^{A \tau} d\tau 
\end{align*}
%
\textbf{Theorem:} $\mathcal{N}[ \mathbf{Q}(t) ] = \mathcal{N}[ \mathbf{O} ] \ \forall t > 0$. 

\textbf{Proof:} \\
\textit{Part 1: Show $\mathcal{N}[ \mathbf{O} ] \subset \mathcal{N}[ \mathbf{Q}(t) ]$ for all $t > 0$.}\\ 
If $x_u \in \mathcal{N}[ \mathbf{O} ]$, then we know from our previous theorem that $C e^{A \tau} x_u = 0$ for all $\tau \ge 0$. 
Let's analyze if $x_u$ is in the null-space of $\mathbf{Q}(t)$:
%
\begin{align*}
    \mathbf{Q}(t) x_u = \int\limits_{0}^{t} e^{A^T \tau} C^T \underbrace{(C e^{A \tau} x_u)}_{0} d \tau
    = 0 \implies x_u \in \mathcal{N}[ \mathbf{Q}(t) ]
\end{align*}
%

\textit{Part 2: Show $\mathcal{N}[ \mathbf{Q}(t) ] \subset \mathcal{N}[ \mathbf{O} ]$ for all $t > 0$.}\\
Let $x_u \in \mathcal{N}[ \mathbf{Q}(t) ]$, which implies $\mathbf{Q}(t) x_u = 0$. Pre-multiplying by $x_u^T$ gives:
\begin{align*}
    x_u^T \mathbf{Q}(t) x_u = 0 \iff \int\limits_{0}^{t} x_u^T e^{A^T \tau} C^T C e^{A \tau} x_u d \tau = \int\limits_{0}^{t} \| C e^{A \tau} x_u \|_2^2 d \tau = 0
\end{align*}
%
Since the integrand is the squared Euclidean norm (which is strictly non-negative) and is continuous, the integral can only evaluate to zero if the integrand is identically zero:
%
\begin{align*}
  C e^{A \tau} x_u = 0 \quad \forall \tau \in [0 , t]
\end{align*}
Taking successive derivatives with respect to $\tau$ and evaluating at $\tau = 0$:
\begin{align*}
  [ C e^{A \tau} x_u]_{\tau = 0} &= 0 \implies C x_u = 0 \\
  \left[\frac{d}{d\tau}( C e^{A \tau} x_u )\right]_{\tau = 0} &= 0 \implies C A x_u = 0 \\
  &\vdots \\
  \left[\frac{d^{n-1}}{d\tau^{n-1}}( C e^{A \tau} x_u )\right]_{\tau = 0} &= 0 \implies C A^{n-1} x_u = 0
\end{align*}
Thus, $\mathbf{O} x_u = 0 \implies x_u \in \mathcal{N}[ \mathbf{O} ]$. We conclude $\mathcal{N}[ \mathbf{Q}(t) ] \subset \mathcal{N}[ \mathbf{O} ]$.

\begin{exmp}
Show that if $\dot{x} = A x$ is asymptotically stable, then the observability Grammian at $t \to \infty$, $Q := \mathbf{Q}_{\infty}$,
satisfies the following Lyapunov equation
%
\begin{align*}
    A^T Q + Q A = - C^T C
\end{align*}
% 
, and this Lyapunov equation has a (unique) positive-definite solution for $Q$,
if and only if, $(A,C)$ is fully observable.
\end{exmp}

\subsection{Further Results in CT-LTI Observability}

The powerful Principle of Duality states that: \textbf{The pair $(A, C)$ is observable if and only if the pair $(A^T, C^T)$ is reachable.}

In view of this duality, we can use our reachability results to immediately derive various conclusions, tests,
standard and canonical forms, etc., for observable and unobservable systems. The reader is highly encouraged to 
derive the following results, theorems, and claims by referring to the dual results in the reachability lecture.

\textbf{Result 1:} The unobservable sub-space, $\mathcal{N} [\mathbf{O}] $ is $A$ invariant, i.e. $x\in \mathcal{N} [\mathbf{O}]  \ \Rightarrow \ A x \in \mathcal{N} [\mathbf{O}] $. 

\textbf{Result 2 (Modal Test for Unobservability):} The pair $(A,C)$ is \textit{unobservable} $\iff$ there exists a right eigenvector $v \neq 0$ of $A$ (where $A v = \lambda v$) such that $C v = 0$. In this case, $\lambda$ is called an unobservable mode.

\textbf{Result 3 (PBH Rank Test for Observability):} The pair $(A,C)$ is \textit{fully observable} $\iff$ the Popov-Belevitch-Hautus (PBH) observability matrix has full column rank for all complex numbers:
\begin{align*}
\mathrm{rank} \left[ \begin{array}{c} \lambda I - A \\ \hline C \end{array} \right] = n , \ \forall \lambda \in \mathbb{C}
\end{align*}

\textbf{Result 4:} Observability is invariant under state/similarity transformation. 

\subsection{Standard Form of Unobservable Systems}

Let $\dot{x} = A x + B u \ , \ y = C x + D u$ be a system that is not fully observable; then it is convenient and practical to choose coordinates (via similarity transformation) to highlight the unobservable and observable ``spaces''. Let 
%
\begin{align*}
   \mathrm{dim}(\bar{\mathcal{O}}) &= \bar{o} 
   \ \& \
   T = \left[ \begin{array}{c|c} T_1 & T_2  \end{array} \right] \ \mathrm{where} \ T_2 \in \mathbb{R}^{n \times \bar{o}} \ , 
   T_1 \in \mathbb{R}^{n \times (n-\bar{o})} \ , 
   \\
   \bar{A} &= T^{-1} A T \ , \ \bar{C} = C T
\end{align*}
%
Let's choose $T_2$ such that $\mathrm{Ra}(T_2) = \mathcal{N}(\mathbf{O}) = \bar{\mathcal{O}}$. In other words, the columns of $T_2$
span the null-space of the observability matrix (the unobservable subspace).
Similarly, let's choose $T_1$ such that the columns of $T$ are linearly independent. Let's analyze the similarity transformation on $A$:
%
\begin{align*}
    A T &=  A \left[ \begin{array}{c|c} T_1 & T_2 \end{array} \right] = T \bar{A} = \left[ \begin{array}{c|c} T_1 & T_2 \end{array} \right]
    \left[ \begin{array}{c|c} A_{11} & A_{12} \\ \hline A_{21} & A_{22}  \end{array} \right]
    \\
    \implies A T_2 &= T_1 A_{12} + T_2 A_{22}
\end{align*}
%
Note that since the unobservable sub-space is $A$-invariant, $A T_2$ must remain in $\mathrm{Ra}(T_2) = \mathcal{N}(\mathbf{O}) = \bar{\mathcal{O}}$. 
Because the columns of $T_1$ are linearly independent from $T_2$, for $T_1 A_{12} + T_2 A_{22}$ to remain entirely in $\mathrm{Ra}(T_2)$, the component along $T_1$ must be zero. Therefore, we must have $A_{12} = \mathbf{0}$.

Now let's analyze the effect of the transformation on the output matrix:
%
\begin{align*}
    \bar{C} &= C T = C \left[ \begin{array}{c|c} T_1 & T_2 \end{array} \right] = \left[ \begin{array}{c|c} C T_1 & C T_2  \end{array} \right] = \left[ \begin{array}{c|c} C_1 & C_2  \end{array} \right]
\end{align*}
%
Because $T_2$ spans the null-space of the observability matrix $\mathbf{O}$, and $C$ is simply the first block-row of $\mathbf{O}$, any vector in $\mathrm{Ra}(T_2)$ is trivially in the null-space of $C$. Thus:
\begin{align*}
    C_2 &= C T_2 = \mathbf{0}
\end{align*}
%
Therefore, the standard unobservable form takes the structure:
%
\begin{align*}
\dot{\bar{x}} &= \left[ \begin{array}{c|c} A_{11} & \mathbf{0} \\ \hline A_{21} & A_{22}  \end{array} \right] \bar{x} + \left[ \begin{array}{c} B_1 \\ B_2 \end{array} \right] u
\\
y &= \left[ \begin{array}{c|c} C_1 & \mathbf{0} \end{array} \right] \bar{x} + D u
\end{align*}
%
Notice that the eigenvalues of $A_{22}$ govern the unobservable modes of the system. These modes are completely ``hidden'' from the output $y(t)$, meaning their dynamic behavior does not affect the measurements at all. We will formalize this decomposition and structure further when we cover the Kalman Decomposition in an upcoming lecture.

\begin{exmp}
\textbf{Problem:} Consider a sparse $3 \times 3$ system that is already in the standard unobservable form:
\begin{align*}
    \dot{x} &= \begin{bmatrix} -1 & 0 & 0 \\ 1 & -2 & 0 \\ 1 & 0 & -3 \end{bmatrix} x \quad , \quad y = \begin{bmatrix} 1 & 0 & 0 \end{bmatrix} x
\end{align*}
(a) Determine the unobservable subspace.\\
(b) Identify the unobservable mode(s).\\
(c) Demonstrate that an initial condition inside the unobservable subspace yields identically zero output despite internal state evolution.

\textbf{Solution:} \\
\textbf{(a) Unobservable Subspace:} By partitioning, we have the observable part $A_{11} = -1$, the coupling block $A_{21} = \begin{bmatrix} 1 \\ 1 \end{bmatrix}$, and the unobservable block $A_{22} = \begin{bmatrix} -2 & 0 \\ 0 & -3 \end{bmatrix}$. The observability matrix is:
\begin{align*}
    \mathbf{O} = \begin{bmatrix} C \\ C A \\ C A^2 \end{bmatrix} = \begin{bmatrix} 1 & 0 & 0 \\ -1 & 0 & 0 \\ 1 & 0 & 0 \end{bmatrix}
\end{align*}
We can clearly see $\mathrm{rank}(\mathbf{O}) = 1 < 3$. The unobservable subspace is the null-space of $\mathbf{O}$:
\begin{align*}
    \bar{\mathcal{O}} = \mathcal{N}[\mathbf{O}] = \mathrm{span}\left\{ \begin{bmatrix} 0 \\ 1 \\ 0 \end{bmatrix} , \begin{bmatrix} 0 \\ 0 \\ 1 \end{bmatrix} \right\}
\end{align*}

\textbf{(b) Unobservable Modes:} The unobservable block $A_{22}$ is diagonal and has eigenvalues $\lambda_2 = -2$ and $\lambda_3 = -3$. The corresponding right eigenvectors for the full matrix $A$ are simply $v_2 = \begin{bmatrix} 0 & 1 & 0 \end{bmatrix}^T$ and $v_3 = \begin{bmatrix} 0 & 0 & 1 \end{bmatrix}^T$. Because $v_2, v_3 \in \mathcal{N}[\mathbf{O}]$, we get $C v_2 = 0$ and $C v_3 = 0$. The modal responses $e^{-2t}$ and $e^{-3t}$ are completely unobservable. 

\textbf{(c) Zero-Output Demonstration:} If we start the system with an initial condition strictly in the unobservable subspace, say $x(0) = \begin{bmatrix} 0 \\ 1 \\ 1 \end{bmatrix}$, we can calculate the state trajectory $x(t)$ without needing to compute the full matrix exponential $e^{At}$. 

Since the first equation is $\dot{x}_1(t) = -x_1(t)$ and we start with $x_1(0) = 0$, we know $x_1(t) = 0$ for all $t \ge 0$. Substituting $x_1(t) = 0$ into the remaining rows completely decouples the ODEs:
\begin{align*}
    \dot{x}_2(t) &= 0 - 2x_2(t) \quad \Rightarrow \quad x_2(t) = e^{-2t} x_2(0) = e^{-2t} \\
    \dot{x}_3(t) &= 0 - 3x_3(t) \quad \Rightarrow \quad x_3(t) = e^{-3t} x_3(0) = e^{-3t}
\end{align*}
The internal state evolution is therefore:
\begin{align*}
    x(t) = \begin{bmatrix} 0 \\ e^{-2t} \\ e^{-3t} \end{bmatrix}
\end{align*}
However, the measured zero-input output remains identically zero:
\begin{align*}
    y(t) = C x(t) = \begin{bmatrix} 1 & 0 & 0 \end{bmatrix} \begin{bmatrix} 0 \\ e^{-2t} \\ e^{-3t} \end{bmatrix} = 0 \quad \text{for all } t \ge 0
\end{align*}
The dynamic behavior of these hidden modal states cannot be seen from the output.
\end{exmp}

\section{Observability in DT-LTI Systems}

\textbf{Definition:} For an LTI discrete-time state-space representation
%
\begin{align*}
  x[k+1] &= A[k] x + B[k] u
\\
  y[k] &= C[k] x + D[k] u
\end{align*}
%
A state $x_u$ is said to be $m$-step unobservable 
if with $x[0] = x_u$ and $\forall \, u[k] \, k \in [0,m)$ we get the same $y[k] \, k \in [0,m)$ as we would with $x[0] = 0$. 

The set, $\bar{\mathcal{O}}_m$, of all $m$-step unobservable states forms a vector space, $\bar{\mathcal{O}}_m \subset 
\mathbb{R}^n$.

Similar to the CT-LTI case, we will only analyze the zero-input response of the system in our 
observability-related derivations. 

If $x_u \in \bar{\mathcal{O}}_m$, then the zero-input response is exactly zero for the first $m$ steps:
%
\begin{align*}
    y[k] = C A^k x_u = 0 \ , \ \forall k \in [0,m)
\end{align*}
%
Expanding this condition step-by-step yields:
%
\begin{align*}
    \begin{array}{c} 
        k = 0 \implies C x_u = 0 \\
        k = 1 \implies C A x_u = 0 \\
        \vdots \\
        k = m-1 \implies C A^{m-1} x_u = 0
    \end{array}
    \quad \Rightarrow \quad
    \underbrace{\begin{bmatrix} C \\ C A \\ \vdots \\ C A^{m-1} \end{bmatrix}}_{\mathbf{O}_m} x_u = 0
\end{align*}
%
This implies that $x_u$ is $m$-step unobservable if and only if
$x_u \in \mathcal{N}[ \mathbf{O}_m ]$.

\textbf{Theorem:} $\mathcal{N}[ \mathbf{O}_i ] = \mathcal{N}[ \mathbf{O}_n ] \subset \mathcal{N}[ \mathbf{O}_m ]$ for $m < n < i$.

\textbf{Corollary:} If $x_u$ is $n$-step unobservable, then it is $i$-step unobservable $\forall i \ge n$.

In this context for DT systems, the unobservable space is defined as $\bar{\mathcal{O}} = \mathcal{N}[ \mathbf{O}_n ]$
and the system is fully observable if and only if 
\begin{align*}
 \bar{\mathcal{O}} = \{0\} \, \iff \, \mathrm{dim}( \mathcal{N}[ \mathbf{O}_n ] ) = 0 \, \iff \,
 \mathrm{rank}[ \mathbf{O}_n ] = n
\end{align*}

\subsection{DT Observability Grammian}

$m$-step observability Grammian is defined as
\begin{align*}
 Q_m = \mathbf{O}_m^T \mathbf{O}_m = \sum\limits_{i=0}^{m-1} (A^i)^T C^T C A^i 
\end{align*}

\textbf{Theorem:} $\mathcal{N}[ \mathbf{O}_m ] = \mathcal{N}[ \mathbf{Q}_m ]$.

Let $x[k+1] = A x[k]$ be asymptotically stable, then $Q = \mathbf{Q}_{\infty}$ is well defined 
and satisfies the Lyapunov equation below
%
\begin{align*}
 A^T Q A - Q = - C^T C    
\end{align*}
%
and this Lyapunov equation has a (unique) positive-definite solution for $Q$,
if and only if, $(A,C)$ is fully observable.

\vspace{6pt}

\textbf{Corollary:} Total Observability Theorem \\
The following statements are mathematically equivalent for a discrete-time system:
\begin{itemize}
 \item The system is fully observable.
 \item The pair $(A,C)$ is observable.
 \item The observability matrix has full rank: $\mathrm{rank}[\mathbf{O}_n] = n$.
 \item \textbf{Modal Test:} $C v \neq 0$ for all right eigenvectors $v$ of $A$ (i.e., $\forall v \neq 0$ such that $A v = \lambda v$).
 \item \textbf{PBH Test:} $\mathrm{rank} \left[ \begin{array}{c} \lambda I - A \\ \hline C \end{array} \right] = n , \ \forall \lambda \in \mathbb{C}$.
\end{itemize}

% **** This ENDS THE EXAMPLES. DON'T DELETE THE FOLLOWING LINE:
\end{document}